\chapter{Tiện ích}

\newpage

\section{Table}

\begin{table}[h!]
    \centering
    \renewcommand{\arraystretch}{1.3}
    \setlength{\tabcolsep}{6pt}
    \begin{tabular}{|p{1cm}|p{2cm}|p{5cm}|p{6cm}|}
        \hline
        \rowcolor[rgb]{0.8, 0.9, 1.0}
        \textbf{STT} & \textbf{MSSV} & \textbf{Họ và tên} & \textbf{Email} \\ \hline
        1 & 24120010 & Trần Văn Tấn Khôi & 24120010@student.hcmus.edu.vn \\ \hline
        2 & 24120028 & Cao Việt Cường & 24120028@student.hcmus.edu.vn \\ \hline
        3 & 24120035 & Lê Thành Đạt & 24120035@student.hcmus.edu.vn \\ \hline
        4 & 24120038 & Nguyễn Phú Đạt & 24120038@student.hcmus.edu.vn \\ \hline
        5 & 24120054 & Mai Phan Nhật Hoàng & 24120054@student.hcmus.edu.vn \\ \hline
    \end{tabular}
    \caption{Thông tin nhóm}
\end{table}

\section{Lists}

\begin{itemize}[topsep=0pt, itemsep=2.5pt, leftmargin=1.5cm]
    \item Lý thuyết: Tìm hiểu các kiến thức lý thuyết liên quan: khái niệm, thành phần, mô hình triển khai, giải pháp kết nối - cấu hình cho các thiết bị: Laptop - Desktop - Smartphone.
    \item Thực hành - demo: Thực hiện cấu hình các thiết bị để xây dựng mạng WLAN thực tế: quay video cách cấu hình các thiết bị:

    \begin{itemize}[label=\ding{111}, leftmargin=1.5cm]
        \item Cấu hình các chức năng router: đặt tên router (SSID) là tên nhóm (\textit{xxCLCxx-Nxx, ví dụ: 19CLC5-N01})
        \item Cấu hình từng loại thiết bị kết nối vô mạng (cần có minh chứng kiểm tra việc kết nối thành công)
        \item Trong video phải có:

        \begin{itemize}[label=\ding{112}, leftmargin=1.5cm]
            \item Thông tin giải thích từng bước cấu hình (thu tiếng trong quá trình làm)
            \item Người thực hiện cấu hình.
        \end{itemize}
        \item Lưu ý: nếu thiết bị không hỗ trợ ad-hoc thì làm trên máy ảo.
    \end{itemize}
\end{itemize}

\newpage

\section{Codebox}

\begin{codebox}[naiveSolution]{Naive Approach}
    \begin{minted}{c++}
    class Player {
        int currentHealth = 100;
        // Player phụ thuộc trực tiếp vào các class cụ thể này
        HealthBarUI healthBar;
        AudioManager audioManager;
        GameManager gameManager;
        public void takeDamage(int amount) {
            this.currentHealth -= amount;
            // Cập nhật UI - Player phải biết hàm update của UI
            healthBar.updateBar(this.currentHealth);
            // Xử lý âm thanh - Player phải biết logic âm thanh
            if (this.currentHealth < 20)
                audioManager.playHeartbeatSound();
            else
                audioManager.playHitSound();
            // Xử lý logic khi máu người chơi không dương
            if (this.currentHealth <= 0)
                gameManager.remove(this);
        }
    };
    \end{minted}
\end{codebox}

\section{Psuedocode}

\begin{algorithm}[H]
\caption{\textproc{nextFrame()}}
\begin{algorithmic}[1]
\Procedure{nextFrame(file)}{}
    \State $SOI \gets \texttt{\{FF D8\}}$
    \State $EOI \gets \texttt{\{FF D9\}}$

    \State $frame\_data \gets \emptyset$

    \State $data \gets$ đọc 2 byte từ file
    \If{$data$ là rỗng}
        \State \Return \texttt{None}
    \EndIf

    \While{$data \neq SOI$}
        \State $pos \gets$ vị trí hiện tại của con trỏ file
        \State đưa con trỏ file về $(pos - 1)$
        \State $data \gets$ đọc 2 byte tiếp theo
        \If{$data$ là rỗng}
            \State \Return \texttt{None}
        \EndIf
    \EndWhile

    \State $frame\_data \gets SOI$
    \State $last\_byte \gets$ rỗng

    \While{True}
        \State $current\_byte \gets$ đọc 1 byte từ file
        \If{$current\_byte$ là rỗng}
            \State \Return \texttt{None}
        \EndIf

        \State nối $current\_byte$ vào $frame\_data$

        % \If{$last\_byte = \texttt{FF}$ \textbf{and} $current\_byte = \texttt{D9}$}
        %     \State tăng bộ đếm frame
        %     \State \Return $frame\_data$
        % \EndIf

        \State $last\_byte \gets current\_byte$
    \EndWhile
\EndProcedure
\end{algorithmic}
\end{algorithm}